\documentclass[acmsmall,12pt,screen,nonacm]{acmart}
\geometry{letterpaper,margin=1in,headheight=16pt,footskip=18pt}
\settopmatter{printacmref=false,printfolios=true}
\setcopyright{none}
\pagestyle{plain}
\usepackage{mathtools}
\newtheorem{theorem}{Theorem}[section]
\newtheorem{lemma}[theorem]{Lemma}
\newtheorem{proposition}[theorem]{Proposition}
\newcommand{\PMC}{\mathcal P}
\newcommand{\Sep}{\mathcal D}
\newcommand{\comp}{\operatorname{comp}}
\newcommand{\poly}{\operatorname{poly}}
\newcommand{\Insert}{\operatorname{Insert}}
\title{Polynomial-Delay Enumeration of All Potential Maximal Cliques}
\author{Zijian Zeng}
\affiliation{\institution{UCSI University}\city{Kuala Lumpur}\country{Malaysia}}
\email{1002266693@ucsiuniversity.edu.my}
\date{}
\begin{document}
\begin{abstract}
A potential maximal clique of a graph is a maximal clique in an
inclusion-minimal chordal completion. We give a deterministic algorithm
enumerating all nonempty potential maximal cliques of every finite simple
undirected graph exactly once, with polynomial delay before the first output,
between consecutive outputs, and from the last output to termination.
The algorithm allows exponential storage and imposes no bound on output size.

The construction builds on the recursive seed and neighborhood framework of
Berg\'e, Limouzy and Schivre. During minimal-separator enumeration, it emits
helper potential maximal cliques whose preimage sets have size at most the
number of vertices. These emissions make separator prefixes chargeable to
actual outputs. Exact suspension at recursive outputs, together with an
injective lift, permits a simultaneous induction on output-rank time and
termination time with an additive polynomial recurrence. Standard buffering
then supplies worst-case delay, including both boundary waits. We specify
the graph-search neighborhood explicitly and verify that it retains the
structural predecessors used in the source framework.
\end{abstract}
\maketitle
\clearpage

\section{The guarantee and its context}

Let $G=(V,E)$ be a finite simple undirected graph, with $n=|V|$ and $m=|E|$.
A \emph{triangulation} is a chordal supergraph on the same vertex set;
it is \emph{minimal} if no proper intermediate supergraph is chordal.
Minimality is by inclusion, rather than by the number of added edges.
A nonempty set is a \emph{potential maximal clique} (PMC) if it is a maximal
clique of some minimal triangulation. Write $\PMC(G)$ for this family.

\begin{theorem}\label{thm:main}
There is one deterministic algorithm and one fixed polynomial $P$ such that,
on every graph $G$, the algorithm enumerates $\PMC(G)$ without repetition.
Time before the first output, between consecutive outputs, and after the
last output until termination is at most $P(n+m+1)$. If $V=\varnothing$,
it terminates within the same bound without output. Exponential space is allowed.
\end{theorem}

PMCs provide an enumeration interface for exact graph algorithms, notably
algorithms based on minimal triangulations and treewidth. Berg\'e, Limouzy
and Schivre~\cite{BLS26}, abbreviated BLS below, give total time
$O(n^4m|\PMC(G)|)$ in Theorem~1 and explicitly leave polynomial delay as
an open challenge in Section~1.4. A total-time bound of this form allows a
long wait early in the computation or after the final output; it does not
by itself prove Theorem~\ref{thm:main}.

The nearest delay guarantees have different scopes. Grigoriev, Kobayashi,
Tamaki and van der Zanden~\cite{GKTZ25}, Theorem~32, give polynomial delay
for triconnected plane graphs. Kenig~\cite{Kenig26}, Theorem~1.1, gives delay
$2^{O(k^3)}n^4\log n$ for PMCs of size at most $k$, including prescribed and
excluded vertices; setting $k=n$ does not give input-polynomial delay.
Brosse, Conte, Limouzy, Punzi and Rucci~\cite{BCLPR24}, Lemma~4, obtain
output-polynomial time with polynomial space. The present algorithm covers
all graphs and all sizes but does not establish polynomial space, an improved
treewidth running-time exponent, or the constrained-search guarantee of
\cite{Kenig26}.

Our contribution is a streaming modification and timing analysis of the BLS
framework. Recursive outputs are processed immediately, and a minimal
separator produces a valid helper PMC before separator enumeration resumes.
A witness bounds each helper fiber by $n$, so an arbitrarily long separator
prefix requires proportionally many already generated solutions. We prove
both a linear bound in every output rank and a linear total bound, with
polynomial input factors. The second invariant pays for recursive finalization
and the final delay. The recurrence adds local costs instead of multiplying
the previous dimension's coefficient.

PMC recognition, recursive lifting and structural predecessor existence are
inherited from BLS and earlier work. Separator extension already appears in
\cite[Proposition~29]{GKTZ25}; the output-count inequality between separators
and PMCs is not claimed here as an innovation. Conversion from polynomial
amortized delay to worst-case delay with a potentially exponential FIFO is
standard; see Capelli and Strozecki~\cite{CS23,CS26}. The additional guarantee
comes from proving the requisite prefix bound for this PMC generator, rather
than applying buffering to an aggregate total-time bound.

\section{Conventions and established primitives}\label{sec:primitives}

For $X\subseteq V$, write $G-X$ for vertex deletion and $N(X)$ for the open
neighborhood outside $X$. When $v\notin X$, let $C_G(v,X)$ be the component
of $G-X$ containing $v$. A component $C$ of $G-S$ is \emph{full} for $S$ if
$N(C)=S$. A set $S$ is a minimal separator precisely when it has at least
two full components. Let $\Sep(G)$ be the set of minimal separators.

We use the following standard facts, with theorem numbers referring to the
specified version of BLS~\cite{BLS26}.
\begin{enumerate}
\item A nonempty $\Omega$ is a PMC if and only if no component of $G-\Omega$
has neighborhood $\Omega$ and every nonadjacent pair of vertices in $\Omega$
belongs to the neighborhood of some component of $G-\Omega$
(Theorem~4 and Lemma~7). This is a polynomial recognition test.
\item For a vertex $a$ and $X\in\PMC(G-a)$, at least one of $X$ and
$X\cup\{a\}$ belongs to $\PMC(G)$ (Lemma~11). In fact the extension is unique
by \cite[Proposition~1]{BCLPR24}; only existence is needed here.
Every PMC containing $a$ is $X\cup\{a\}$ for a child PMC $X$, or
$S\cup\{a\}$ for $S\in\Sep(G)$ (Lemma~12).
\item Minimal separators admit polynomial-delay enumeration, as stated
in BLS Lemma~1 from Berry, Bordat and Cogis~\cite{BBC99}.
\item A minimal triangulation can be computed in polynomial time by MCS-M
of Berry, Blair, Heggernes and Peyton~\cite{BBHP04}, Theorem~4.4.
\end{enumerate}

For clarity about endpoints in the separator primitive, one may implement
the known component-neighborhood scheme with a queue and a trie. Initialize
with $N(C)$ for all components $C$ of $G-N[v]$, for every vertex $v$.
On dequeuing $S$, emit it; then, for every $x\in S$, insert $N(C)$ for
components $C$ of $G-(S\cup N(x))$. Insert only unseen sets and, if desired,
validate the two-full-components criterion. The known completeness theorem
supplies all separators. Initialization and each queue expansion require
polynomially many graph searches. Emitting \emph{before} expansion gives
polynomial initial, successive and final delay, including the empty-family
case. No exhaustive subset oracle is part of this primitive.

For MCS-M, the path tests in its primary pseudocode can be performed naively
by graph search in the fixed input block. All tests for a numbering step use
the weights before that step's increments. There are polynomially many tests,
so no optimized implementation is required here. A perfect elimination order
of the resulting chordal graph lists its maximal cliques in polynomial time:
take each vertex with its later neighbors and discard nonmaximal candidates.

Sets are $n$-bit vectors. Use binary tries of depth at most $n$ for dictionaries
and linked queues, so no operation scans all previous outputs. Graph searches,
recognition and serialization have polynomial worst-case cost. Store at most
exponentially many records at each of at most $n$ recursive levels; pointers
and counters therefore have polynomial bit length. We use a pointer RAM with
polynomial-length words, charging polynomial bit cost to its operations.
Memory may be retained until termination: an exponential deallocation pass
must not be appended after the last output.

\section{A separator helper with a witnessed fiber bound}\label{sec:helper}

\begin{lemma}\label{lem:helper}
For a connected graph $G$, a deterministic polynomial-time map
$\Gamma:\Sep(G)\to\PMC(G)$ can be chosen so that every image has at most
$n$ preimages. More precisely, for each $S$, some component $D$ of
$G-\Gamma(S)$ satisfies $N_G(D)=S$.
\end{lemma}

\begin{proof}
Choose two full components $C,D$ of $G-S$ using fixed label order. For every
component $C_i$ of $G-S$, make $S$ a clique in $G[S\cup C_i]$ and compute a
minimal triangulation $T_i$ of this saturated block. Their union $T$ is
chordal: its pieces are chordal, intersect in the clique $S$, and have no
edges between different $C_i$.

We check minimality with respect to the original graph. Let $H$ be chordal
with $G\subseteq H\subseteq T$. If an edge $xy$ inside $S$ is absent from
$H$, fullness and connectivity give an $x$--$y$ path with interior in $C$
and another with interior in $D$. Choose shortest such paths in $H$ on their
respective sides. Both are induced and have length at least two. Their
interiors have no cross edges in $T$ and hence none in $H$, so their union
is a chordless cycle of length at least four. This contradicts chordality.
Thus $S$ is a clique in $H$. Each $H[S\cup C_i]$ is now a chordal supergraph
of the saturated block contained in $T_i$, and equals $T_i$ by minimality.
Consequently $H=T$, proving that $T$ is a minimal triangulation of $G$.

There is a maximal clique of $T_C$ properly containing $S$. To justify this
extension, take a clique tree of the connected chordal graph $T_C$.
For each $s\in S$, the nodes whose cliques contain $s$ form a subtree;
these subtrees pairwise intersect because $S$ is a clique. The union of nodes
whose cliques meet $C$ is also a connected subtree: the subtrees for vertices
of connected $C$ are connected through its edges. It meets each subtree for
$s\in S$ because $C$ is full. The Helly property for subtrees of a tree gives
a node containing all of $S$ and a vertex of $C$.
For completeness, tree Helly follows by rooting the tree and choosing the
deepest of the nearest-to-root vertices of the pairwise intersecting subtrees;
that vertex lies in every subtree.

Choose deterministically a maximal clique $\Gamma$ of $T_C$ containing $S$
properly. It is maximal in $T$ too, since no other block contains a vertex
adjacent to its vertex in $C$. Thus $\Gamma\in\PMC(G)$.
We have $S\subsetneq\Gamma\subseteq S\cup C$, so $D$ survives as a whole
component of $G-\Gamma$ and $N_G(D)=S$. For a fixed image $\Gamma$, distinct
preimage separators are distinct neighborhoods of components of $G-\Gamma$.
There are at most $n$ such components. All choices and computations above
take polynomial time.
\end{proof}

The useful assertion is about the \emph{actual map}: after processing any
$r$ distinct separators, at least $r/n$ distinct helper PMCs have appeared.
They may coincide with earlier outputs of another phase; those earlier outputs
still pay for their separator preimages. A final inequality
$|\Sep(G)|\le n|\PMC(G)|$ alone would not account for a prefix if the
separator phase did not actually emit these helpers.

\section{An explicit neighborhood and the inherited preimages}\label{sec:arcs}

Call a PMC $\Omega$ \emph{$a$-nice} if $a\notin\Omega$ and
$\Omega\notin\PMC(G-a)$, following BLS Definition~16. Non-$a$-nice PMCs
are the seed family. BLS Theorems~8--10 supply predecessors for $a$-nice
targets. We state the portion of their structural construction used here,
and then connect it to an explicit executable neighborhood.

\begin{proposition}[Structural witnesses from BLS]\label{prop:witness}
Let $\Omega$ be $a$-nice, let $C_a=C_G(a,\Omega)$, $S_0=N(C_a)$ and
$T_0=\Omega\setminus S_0$. Then $T_0$ is nonempty and lies in a full
component $C_T\ne C_a$ of $G-S_0$. There are $W\in\PMC(G)$ and
$u\in W\cap C_a$ with $T_0\subseteq W$ and
\begin{equation}\label{eq:sandwich}
 \Omega\subseteq X:=W\cup N(u)\subseteq\Omega\cup C_a.
\end{equation}
Moreover, either $a\in W$, or $a\notin W$ and
$C_G(a,W)\subsetneq C_a$.
\end{proposition}

These are inherited structural facts, rather than a new characterization.
The placement of $T_0$ is BLS Lemma~10. To make the provenance of
\eqref{eq:sandwich} explicit, Theorem~9 constructs
$W=\{a\}\cup\widehat S\cup T_0$ with $u=a$, and every missing vertex
of $S_0$ is adjacent to $a$. Theorem~10 constructs
$W=D\cup(\Omega\setminus S_D)\cup S'_D$, where $D\subseteq C_a$,
$S'_D\subseteq S_D$, and every missing vertex of $S_D$ is adjacent to the
selected $u\in D$. These statements give both inclusions in
\eqref{eq:sandwich} and retain $T_0$. The latter theorem also gives the
strict decrease of the $a$-component. Existence of the requisite alternatives
is supplied by Theorem~8.

Here is the neighborhood actually computed. For $W\in\PMC(G)$ and each
ordered pair of distinct $u,b\in W$:
\begin{enumerate}
\item Put $X=W\cup N(u)$ and let $A=C_G(b,X\setminus\{b\})$.
\item Reject if $u\in A\cup N(A)$. Otherwise let
$U=C_G(u,N(A))$ and $S=N(U)$.
\item Let $B=C_G(b,S)$ and $Z=X\cap(S\cup B)$. Retain $Z$ exactly when
it is a nonempty PMC.
\end{enumerate}
All component queries are defined. In particular, $b\in A$ and
$S\subseteq N(A)$, so $b\notin S$; the guard ensures $u\notin N(A)$.
There are at most $n^2$ pairs and polynomial work per pair.

BLS Definition~12 assumes that $X\setminus\{a\}$ is an $(a,b)$-separator,
whereas Definition~20 chooses both endpoints in $X$. We avoid relying on this
endpoint wording by defining the graph searches above. The following argument
shows directly that they keep every structural arc needed for reachability.

\begin{lemma}\label{lem:transfer}
For each target and witness in Proposition~\ref{prop:witness}, every choice
$b\in T_0$ makes the pair $(u,b)$ survive the guard and return $Z=\Omega$.
Consequently all PMCs are reachable from the non-$a$-nice PMCs by this
neighborhood.
\end{lemma}

\begin{proof}
Since $S_0\subseteq X\setminus\{b\}$, the component $A$ stays in $C_T$.
It has no vertex or neighbor in $C_a$, a different component of $G-S_0$.
For any $s\in S_0$, the PMC recognition criterion supplies either the edge
$bs$ or a $b$--$s$ path with interior in one component of $G-\Omega$.
In the latter case the interior lies in $C_T$, since the path cannot leave
that component of $G-S_0$ before reaching $s$. It therefore avoids $X$ by
\eqref{eq:sandwich}. Starting at $b$, the search defining $A$ reaches a
neighbor of $s$. Thus
\[
 S_0\subseteq N(A),\qquad N(A)\cap C_a=\varnothing.
\]
The guard keeps $u$. Removing $N(A)$ isolates the connected set $C_a$ from
everything else because it removes all of $S_0=N(C_a)$ and none of $C_a$.
Hence $U=C_a$ and $S=S_0$. Now $B=C_T$, and
\eqref{eq:sandwich} gives $X\cap(S_0\cup C_T)=\Omega$. Recognition accepts it.

Follow structural predecessors backwards from an $a$-nice PMC. Until a
non-$a$-nice predecessor is reached, its nonempty $a$-component strictly
decreases. This process cannot continue indefinitely and takes at most $n$
steps. Reverse the sequence and use the proved computational arcs.
\end{proof}

\section{The streaming generator}\label{sec:generator}

For now assume that $H$ is connected and nonempty. Define a recursive generator
$\mathcal A(H)$ which may suspend at an emission and resume from exactly that
point. If $H$ is complete, emit $V(H)$ and halt. Otherwise choose $a$ so that
$H-a$ is connected, for example a leaf of a fixed spanning tree.

Maintain an empty trie $\mathrm{Seen}$ and an empty linked queue $L$.
The operation $\Insert(Z)$ checks the trie; if $Z$ is unseen, it stores $Z$,
appends $Z$ to $L$ and immediately emits $Z$, suspending before further work.
All calls below pass recognized PMCs.
\begin{enumerate}
\item[\textbf{I.}] Run $\mathcal A(H-a)$ up to its next emission $X$ and
suspend the child. Test $X$ and $X\cup\{a\}$ in fixed order for membership
in $\PMC(H)$, inserting each valid set. Resume the child only after these
tests and their emissions are finished. Repeat until the child halts.
\item[\textbf{II.}] Run the polynomial-delay minimal-separator enumerator
on $H$. Suspend it at each distinct separator $S$. Compute and insert
$\Gamma(S)$ from Lemma~\ref{lem:helper}; then test $S\cup\{a\}$ and insert
it if valid. Resume separator enumeration afterwards, continuing to its halt.
\item[\textbf{III.}] While $L$ is nonempty, pop its first entry $W$,
compute the neighbors of Section~\ref{sec:arcs} in fixed order, and insert
each valid neighbor. Halt when the queue is empty.
\end{enumerate}
The queue is filled at discovery, including in phases I and II. No completed
exponential seed list is copied at a phase transition.

\begin{lemma}\label{lem:correct}
The generator terminates and emits exactly $\PMC(H)$ without repetition.
Different child PMCs have different parent lifts, so
$|\PMC(H-a)|\le |\PMC(H)|$.
\end{lemma}

\begin{proof}
Recognition and the trie give validity and uniqueness. Each child PMC has at
least one valid lift, and deleting $a$ from a lift recovers the child set.
Thus distinct children supply distinct parent outputs. Phase I includes
all parent PMCs that are PMCs of $H-a$. Every PMC containing $a$ has one
of the forms supplied in Section~\ref{sec:primitives}, and is found in
phase I or II. These are exactly the non-$a$-nice PMCs. Extra helper seeds
do not reduce reachability, so Lemma~\ref{lem:transfer} proves completeness
of phase III. At most $2^{|V(H)|}$ sets are inserted, each queue entry is
popped once, and all phases consist of finite operations; recursion reduces
the vertex count. Termination follows.
\end{proof}

\section{One induction for output ranks and termination}\label{sec:timing}

Let $K=|\PMC(H)|\ge1$, $K_c=|\PMC(H-a)|$ and $D=|\Sep(H)|$.
We measure generator work, rather than time while a coroutine is externally
suspended. Choose a fixed nondecreasing polynomial $p(n)\ge1$ bounding each
of: initialization and a complete-graph case; local processing of one child
output; one separator gap together with its helper and lift processing; one
entire queue expansion; and each of the boundedly many phase transitions
and finalization operations. Include dictionaries, queues, serialization,
coroutine bookkeeping and their polynomial bit costs. Such a computable
polynomial exists by the specified implementations. Fix the implementation
once, so its constants are independent of the input and output counts.
Take $Q(1)$ large enough and set
\begin{equation}\label{eq:recurrence}
 Q(n)=Q(n-1)+(n+10)p(n)\quad(n\ge2).
\end{equation}
For instance $Q(1)+(n+1)(n+10)p(n)$ is a fixed polynomial upper bound.

\begin{lemma}\label{lem:timing}
For every connected nonempty $n$-vertex graph, simultaneously:
\begin{align}
 E_H(j)&\le Q(n)j &&(1\le j\le K),\label{eq:prefix}\\
 T_H&\le Q(n)K,\label{eq:termination}
\end{align}
where $E_H(j)$ is the work through the $j$-th distinct emission and $T_H$ is
the work through complete termination.
\end{lemma}

\begin{proof}
Induct on $n$, proving both assertions together. Complete graphs obey the
bounds by the choice of $p$ and $Q$. For a noncomplete graph,
Lemma~\ref{lem:correct} gives $K_c\le K$, Lemma~\ref{lem:helper} gives
$D\le nK$, and at most $K$ queue entries are expanded. By the child's
termination bound, its total cost is at most $Q(n-1)K_c$. All local work
is at most
\[
 p(n)(K_c+D+K+5)\le(n+7)p(n)K.
\]
Together with \eqref{eq:recurrence} this proves \eqref{eq:termination}.
In particular the recursive coefficient is used once, not multiplied by an
additional constant greater than one.

For \eqref{eq:prefix}, classify the $j$-th emission by its phase.

\emph{Phase I.} Suppose it occurs while processing child emission $r$.
The first $r-1$ child outputs already supplied distinct parent lifts, so
$r\le j$. Because the child is suspended at emission $r$, its elapsed work
is at most $Q(n-1)r$ by the child prefix bound. Local work so far is at most
$p(n)(r+2)\le3p(n)j$, proving the desired estimate. An emission during
the first lift test does not resume the child before the second test.

\emph{Phase II.} The child has halted and all its lifts are present, so
$K_c\le j$. If the current separator has rank $r$, its $r-1$ predecessors
have all had their helper PMCs emitted earlier, either in this phase or
already in phase I. Each image has at most $n$ preimages, hence
\[
 r-1\le n(j-1),\qquad r\le nj.
\]
The child cost is at most $Q(n-1)K_c$ by its termination bound. Local cost
is at most $p(n)(K_c+r+4)\le(n+5)p(n)j$. These estimates prove
\eqref{eq:prefix}, also when the current output is the separator lift rather
than a newly emitted helper.

\emph{Phase III.} Every separator helper has been emitted before this phase,
so $D\le n(j-1)$. If $q$ queue expansions have begun by this emission,
each popped set was emitted earlier and is popped once; thus $q\le j$.
Also $K_c\le j$. Child and local costs are at most
\[
 Q(n-1)K_c+p(n)(K_c+D+q+5)
 \le [Q(n-1)+(n+7)p(n)]j.
\]
This is bounded by $Q(n)j$. The simultaneous induction is complete.
\end{proof}

The distinction between the two invariants is essential. A recursive prefix
bound alone does not control the child's final tail before phase II.
The termination bound pays for that tail with the child's already lifted
outputs. Conversely a total-time bound alone would not control the number
of processed separators before the next fresh output. Their actual helper
emissions provide that control.

\section{Buffering and all input boundaries}\label{sec:buffer}

We give the standard FIFO conversion explicitly to include termination.
For a connected nonempty graph, use an integer polynomial upper bound
$q\ge Q(n)$. Simulate $\mathcal A(H)$ in blocks of $q$ steps, stopping a
block early if the generator halts, and put its emissions in a FIFO buffer.
After each block return one buffered solution. Once simulation has halted,
drain the buffer with polynomial serialization time per solution and halt.
Account for the polynomial enqueue cost in an emission's simulated cost;
step counting and simulation overhead change the polynomial bound only by a
polynomial factor.

At the end of block $j\le K$, the prefix bound ensures at least $j$
generated solutions, whereas only $j-1$ have been returned. Thus the buffer
is nonempty. If $j>K$, the termination bound says the generator has already
halted by $qK<qj$, so no empty-buffer block can occur with it still running.
If it halts within a block, all unreturned solutions are already buffered.
It halts no later than the end of block $K$. Consequently after the $K$-th
returned solution there is no remaining simulation tail. The first wait,
each intermediate wait, and the final wait are polynomial. This is the
usual amortized-to-delay buffering principle~\cite{CS23,CS26}, applied to
the two bounds just proved; it makes no polynomial-space assertion.

It remains to include disconnected and empty graphs. In a minimal
triangulation of a disconnected graph there are no edges between original
connected components: deleting all such edges leaves the disjoint union
of chordal induced subgraphs containing the original graph, so minimality
forbids their presence. Conversely, a disjoint union of minimal triangulations
of its components is minimal, since a proper chordal subtriangulation would
restrict to a proper one on some component. Thus the nonempty PMCs of $G$
are exactly the PMCs of its nonempty connected components, on their original
labels.

Find these components in polynomial time. The zero-vertex graph halts without
output. Otherwise run the buffered connected algorithm on the components
sequentially. Each has an output, and nonempty sets from different components
are distinct. At a component boundary, the previous final wait plus the next
initial wait is polynomial in $n+m+1$. Isolated vertices and complete
components are included in the base case. This proves Theorem~\ref{thm:main}.

\section{Verification and scope}

The proof uses established structural theorems as specified above; finite
experiments are supplementary diagnostics. A retained test record compares
complete PMC and separator families on 29 specified graphs of order at most
ten, with 471 PMC outputs in total, against exhaustive finite oracles.
It also checks helper witnesses, fiber sizes, lifts, queue counts and absence
of repeated outputs. A separate automated check compares the literal
minimal-triangulation definition with the PMC criterion on all 1100 labelled
graphs through order five, including the empty graph. It also independently
checks the 29 serialized output families using that criterion. These
checks do not prove the universal theorem or its asymptotic timing bounds.
Exhaustive subsets are used by the finite verification oracles and are absent
from the theoretical enumeration algorithm.

The theorem guarantees progress for every requested output prefix. It does
not bound the number of PMCs polynomially, promise polynomial memory, rank
outputs, or provide the inclusion/exclusion oracle of \cite{Kenig26}.
The polynomial exponent is not optimized. Practical implementations and a
polynomial-space delay algorithm require further work. The literature
comparison concerns the cited primary versions; it does not certify
exhaustive priority.

Generative AI materially assisted development of the proof and computational
checks, as detailed in the statement at the end. Neither a model review nor
a successful finite check is a formal proof certificate.

\clearpage
\AtBeginEnvironment{thebibliography}{\interlinepenalty=10000}
\bibliographystyle{alpha}
\bibliography{references}

\section*{AI methodology statement}
OpenAI Codex, based on GPT-6, was used substantially for mathematical research:
developing and attacking the separator helper, the explicit neighborhood
transfer, the recursive streaming algorithm and its timing induction;
writing and running finite diagnostic code; consulting primary literature;
and drafting and building this manuscript. A separate Codex context performed
an internal mathematical and novelty review of the frozen argument and
additional finite verification. Such review is neither external peer review
nor formal verification, and no independent human verification of all
mathematical claims or references is asserted here. The proof is presented
in full with its source dependencies; the listed human author is responsible
for assessing its accuracy, originality and suitability before any submission.
\end{document}
